\abstract

\keywords{ИНФОРМАЦИОННАЯ СИСТЕМА, МОЛОЧНАЯ ФЕРМА, RUST, SVELTEKIT, МАШИННОЕ ОБУЧЕНИЕ, ПРЕДИКТИВНАЯ АНАЛИТИКА, ПРОГНОЗИРОВАНИЕ НАДОЕВ, ОБНАРУЖЕНИЕ МАСТИТА, ИНТЕГРАЦИЯ LELY, DOCKER}

Объектом исследования является процесс автоматизации управления молочной фермой.

Предметом исследования --- методы и средства разработки информационных систем управления молочными фермами с интеграцией роботизированного доильного оборудования и предиктивной аналитикой на основе моделей машинного обучения.

Цель работы --- разработка информационной системы, обеспечивающей комплексный учёт животных, надоев молока, воспроизводства, кормления, фитнес-мониторинга, пастьбы и интеграцию с роботизированным доильным оборудованием Lely, а также модуля предиктивной аналитики на основе моделей машинного обучения.

Основное содержание работы состоит в проектировании и реализации трёхзвенной архитектуры: серверной части на языке Rust с фреймворком Axum (25 модулей бизнес-логики, JWT-аутентификация, middleware безопасности, интеграция с Lely Horizon), клиентского веб-приложения на SvelteKit (более 24 страниц, адаптивный интерфейс, более 16 UI-компонентов) и модуля аналитики с 11 моделями машинного обучения на Python/FastAPI (XGBoost, Prophet, KMeans, Isolation Forest). Серверная часть содержит собственные алгоритмы прогнозирования: метод Хольта--Винтерса, кривую лактации Вуда, индекс здоровья на основе Z-оценок и систему сравнения 8 моделей прогнозирования временных рядов с автоматическим выбором оптимальной. База данных PostgreSQL (более 35 таблиц) использует TimescaleDB для эффективной работы с временными рядами.

Основным результатом работы является информационная система управления молочной фермой, покрывающая полный цикл управления стадом: от учёта и мониторинга до предиктивной аналитики с доверительными интервалами и раннего обнаружения заболеваний (мастит, кетоз, охота).

Данные результаты предназначены для автоматизации процессов учёта и мониторинга на молочных фермах, оснащённых роботизированным доильным оборудованием, повышения эффективности управления стадом и снижения экономических потерь от заболеваний животных.

Внедрение системы позволяет сократить время на ведение документации, повысить точность прогнозирования продуктивности стада и своевременно выявлять заболевания на ранних стадиях, когда клинические признаки ещё не очевидны.
